\begin{corollary}\label{tmh:c2}  Sean $A\subseteq\Rn{n}$ abierto y arcoconexo y $\mathbf{F}:A\subseteq\Rn{n}\to\Rn{n}$  un campo vectorial conservativo entonces para todo par de curvas $C_1$ , $C_2$ simples en $A$, con los mismos puntos extremos y misma orientaci\'on, se tiene que
    \[
    \int_{C_1}\mathbf{F}\cdot d\mathbf{s}=\int_{C_2}\mathbf{F}\cdot d\mathbf{s}. 
    \]
    Dada esta igualdad, se dice que la integral de línea de un campo conservativo tiene \textit{independencia del camino}.

    \textit{Demostración.}   Sean $\boldsymbol{\sigma}_1:[a_1,b_1]\to A$  y $\boldsymbol{\sigma}_2:[a_2,b_2]\to A$  parametrizaciones de  $C_1$ y $C_2$ respectivamente que preservan sus orientaciones y sea  $f$ un funci\'on potencial de $\mathbf{F}$. Por definici\'on  (\ref{int_cv}) se tiene que:
    \begin{align*}
        \int_{C_1}\mathbf{F}\cdot d\mathbf{s}&=\int_{\boldsymbol{\sigma}_1}\grad f\cdot d\mathbf{s}\\
        \int_{C_2}\mathbf{F}\cdot d\mathbf{s}&=\int_{\boldsymbol{\sigma}_2}\grad f\cdot d\mathbf{s}
    \end{align*}
    Además, como $C_1$ y $C_2$ tienen lo mismos puntos extremos, entonces $\boldsymbol{\sigma}_1(a_1)=\boldsymbol{\sigma}_2(a_2)\;\land\;\boldsymbol{\sigma}_1(b_1)=\boldsymbol{\sigma}_2(b_2)$.
    Por lo tanto,  por el Corolario \ref{tmh:c1}:
    $$  \int_{\boldsymbol{\sigma}_1}\grad f\cdot d\mathbf{s}=\int_{\boldsymbol{\sigma}_2}\grad f\cdot d\mathbf{s} $$
     lo que demuestra que:   $$ \int_{C_1}\mathbf{F}\cdot d\mathbf{s}=\int_{C_2}\mathbf{F}\cdot d\mathbf{s}$$
   
\end{corollary}




\begin{corollary}\label{tmh:c2} 
Sean $A \subseteq \mathbb{R}^n$ abierto y arcoconexo y $\mathbf{F}: A \subseteq \mathbb{R}^n \to \mathbb{R}^n$ un campo vectorial conservativo. Entonces, para todo par de curvas $C_1$ y $C_2$ simples en $A$, con los mismos puntos extremos y misma orientación, se tiene que: 
\[ \int_{C_1} \mathbf{F} \cdot d\mathbf{s} = \int_{C_2} \mathbf{F} \cdot d\mathbf{s} \] 
Dada esta igualdad, se dice que la integral de línea de un campo conservativo tiene \emph{independencia del camino}. 
\end{corollary}

\begin{proof}
Sean $\boldsymbol{\sigma}_1: [a_1,b_1] \to A$ y $\boldsymbol{\sigma}_2: [a_2,b_2] \to A$ parametrizaciones de $C_1$ y $C_2$ respectivamente que preservan sus orientaciones, y sea $f$ una función potencial de $\mathbf{F}$. 

Por la Definición~\ref{int_cv} y sabiendo que $\mathbf{F} = \nabla f$, se tiene que: 
\begin{align*} 
\int_{C_1} \mathbf{F} \cdot d\mathbf{s} &= \int_{\boldsymbol{\sigma}_1} \nabla f \cdot d\mathbf{s} \\ 
\int_{C_2} \mathbf{F} \cdot d\mathbf{s} &= \int_{\boldsymbol{\sigma}_2} \nabla f \cdot d\mathbf{s} 
\end{align*} 

Además, dado que $C_1$ y $C_2$ comparten los mismos puntos extremos y el mismo sentido de recorrido, sus parametrizaciones satisfacen que $\boldsymbol{\sigma}_1(a_1) = \boldsymbol{\sigma}_2(a_2) \ \land \ \boldsymbol{\sigma}_1(b_1) = \boldsymbol{\sigma}_2(b_2)$. 

Por lo tanto, aplicando el Corolario~\ref{thm:c1}, obtenemos: 
\[ \int_{\boldsymbol{\sigma}_1} \nabla f \cdot d\mathbf{s} = \int_{\boldsymbol{\sigma}_2} \nabla f \cdot d\mathbf{s} \] 
Lo que demuestra finalmente:
\[ \int_{C_1} \mathbf{F} \cdot d\mathbf{s} = \int_{C_2} \mathbf{F} \cdot d\mathbf{s} \qedhere \] 
\end{proof}


\begin{corollary}\label{tmh:c3} 
Sean $A \subseteq \mathbb{R}^n$ abierto y arcoconexo y $\mathbf{F}: A \subseteq \mathbb{R}^n \to \mathbb{R}^n$ un campo vectorial conservativo. Entonces, para toda curva $C$ simple cerrada contenida en $A$, se cumple que: 
\[ \oint_C \mathbf{F} \cdot d\mathbf{s} = 0 \] 
\end{corollary}

\begin{proof}
Sea $f$ una función potencial de $\mathbf{F}$ y sea $\boldsymbol{\sigma}: [a,b] \to A$ una parametrización de $C$ que preserva su orientación. Entonces, por la Definición~\ref{int_cv} y sabiendo que $\mathbf{F} = \nabla f$, tenemos: 
\begin{equation*} 
\oint_C \mathbf{F} \cdot d\mathbf{s} = \int_{\boldsymbol{\sigma}} \nabla f \cdot d\mathbf{s} 
\end{equation*} 

Por el Teorema~\ref{thm:t1}, sabemos que: 
\begin{equation*} 
\int_{\boldsymbol{\sigma}} \nabla f \cdot d\mathbf{s} = f(\boldsymbol{\sigma}(b)) - f(\boldsymbol{\sigma}(a)) 
\end{equation*} 

Además, dado que la curva $C$ es cerrada, su punto final coincide con su punto inicial, lo que significa que $\boldsymbol{\sigma}(b) = \boldsymbol{\sigma}(a)$. Por la buena definición de la función $f$, se cumple que $f(\boldsymbol{\sigma}(b)) = f(\boldsymbol{\sigma}(a))$. 

Sustituyendo esta igualdad en la ecuación anterior, obtenemos:
\[ f(\boldsymbol{\sigma}(b)) - f(\boldsymbol{\sigma}(a)) = 0 \]
Lo que demuestra finalmente que: 
\[ \oint_C \mathbf{F} \cdot d\mathbf{s} = 0 \qedhere \] 
\end{proof}








